% ============================================================================
% A 题 · 论文 LaTeX 前言（preamble）—— 由 50\_论文/02\_章节稿/gen\_latex.py 生成
% GEN\_LATEX\_VERSION: gen\_latex/v11
% 设计依据：kb《AI论文插图与LaTeX排版》＋ 官方《论文格式规范》＋ 我方《内容与排版绘图交接说明》§四
% 编译：xelatex 两遍（见 README\_排版交付说明.md）
% ============================================================================
% ★ 行距用 **ctexart 的类选项**指定（`linespread=<值>`），**不可用 `\linespread`**：
%   ctex 会在 \begin{document} 处按自身方案重设 \baselinestretch，故前言区的 `\linespread{}` 会被覆盖
%   （实测：`\linespread{1.12}` 与 `\linespread{1.09}` 页数完全相同，证实其为**空操作**）。
\documentclass[12pt,a4paper,linespread=1.0]{ctexart}

\usepackage{geometry}\geometry{left=2.5cm,right=2.5cm,top=2.5cm,bottom=2.5cm}  % 官方第一条

% ---- ★ v4 新增：版面调参区（属**模板接口**；贵方如需自定义，改本区块即可，或取回 preamble.tex.bak）----
% 依据：官方第八条明示"论文中的**字号、字体、行距**不作统一要求"，故以下取值均为可选优化，非官方要求。
% 目的：正文页数由"内容侧已触底"后的 32 页收进官方第四条红线（≤30 页），并消除版面松散感。
% 行距：见上方 \documentclass 的 `linespread=1.0` 类选项（前置件的 `\linespread` 会被 ctex 覆盖）。
%   ★ 官方依据：《全国大学生数学建模竞赛论文格式规范》**第八条**——"字号、字体、行距、颜色等
%   **不做统一要求**" ⟹ 行距可自主设定。由 1.06 调至 **1.0（单倍）**以压缩页数。
%   **注意不可再动页边距**：**第一条**要求"上下左右各留出**至少 2.5 厘米**"，本件
%   `\geometry{...=2.5cm}` 恰好压在下限，**再小即违规**。
% 浮动体密度：允许每页承载更多浮动体，避免"整页只有一张图/表"的松散页
\setcounter{topnumber}{4}
\setcounter{bottomnumber}{2}
\setcounter{totalnumber}{6}
\renewcommand{\topfraction}{0.92}
\renewcommand{\bottomfraction}{0.90}
\renewcommand{\textfraction}{0.06}
\renewcommand{\floatpagefraction}{0.75}
% -----------------------------------------------------------------------------
\usepackage{amsmath,amssymb}
\usepackage{booktabs}          % 三线表（\toprule/\midrule/\bottomrule），全篇无竖线
\usepackage{longtable}         % 长表：自动重复表头（对应"（续）"要求）
\usepackage{array}
% 图片搜索路径：**相对路径实时链接** 50\_论文/05\_成品图（禁用绝对路径），并保留本地 figures/ 备选
\usepackage{graphicx}\graphicspath{{../../05\_成品图/}{figures/}}
\usepackage{float}
\usepackage{wrapfig}           % ★ v14：个别图改"环绕"排版（用户指令；正文绕图，图不独占版面）
\usepackage[font=small,labelsep=quad]{caption}
\usepackage{listings}
% ★ v12：**新字符兜底**（`literate` 在 XeLaTeX 下对多字节字符不总生效，故再叠加 `newunicodechar`）——
%   把代码注释里 DejaVu Sans Mono 缺失的两个字形映射到**已有字体**里的同义字形：
%   `⟹`(U+27F9) → 数学体 `\Longrightarrow`；`✅`(U+2705) → `pifont` 的 `\ding{51}`（✓）。
%   两者**都不改动 `code/` 源文件字节** ⟹ 与真源逐字节一致、哈希不变。
\usepackage{newunicodechar}
\newunicodechar{⟹}{\ensuremath{\Longrightarrow}}
\newunicodechar{✅}{\ding{51}}
\usepackage{pifont}            % ★ v12：`\ding{51}` 用于附录 B 代码注释里的 ✅ 字形（见下 literate）
\usepackage{xcolor}
\usepackage[hidelinks]{hyperref}

% ★ v12：**图表"留白"收紧**（用户反馈"很多图表都有很大的白边"；官方第八条"字号、字体、行距、
%   颜色等不做统一要求" ⟹ 可自主设定）。取证：① **图件本体无白边**——19 张成品图实测
%   "页面盒≈内容盒"，四边白边率 0–2.6%（0 张 >12%）；② **页面无大空白**——全篇 126 页中
%   仅 p1（标题页）与 p38（页尾 95pt）空隙 >60pt，其余最大空隙 ≤31pt、填充率 ~100%。
%   ⟹ 观感上的"白边"来自**插入图表时的默认间距**（LaTeX 默认 `\textfloatsep`≈12pt、
%   `\intextsep`≈12pt、`\abovecaptionskip`≈10pt 偏松）。此处统一收紧（中等幅度，不挤读感）。
% ★ v12 归因结论（实测）：把 `\textfloatsep`/`\intextsep`/`\floatsep` 收紧到 8–9pt
%   会改变分页、进而使两张**长表**（`longtable`）的页内分块超出 ⟹ 冒 2 处
%   `Overfull \vbox`（254／260 pt）；**停用即归零**。依"无证据不留改动"纪律，
%   **保留 LaTeX 默认的浮动体间距**，只收紧**题注间距**（下图，实测 0 vbox／0 hbox）。
\setlength{\abovecaptionskip}{5pt}                   % 图/表 ↔ 题注（默认 ≈10pt）
\setlength{\belowcaptionskip}{0pt}
% 注：**未动 `\arraystretch`**（实测把它调紧会与长表 `longtable` 的分块配合出现
%     `Overfull \vbox`，见交付件变更留痕 ⟹ 依"无证据不留改动"回退，保持默认 1.0）。

% ---- 源程序（附录 B）：文件放入 code/ 即自动嵌入，缺位时显示占位框（保证可编译）----
% 注：文件名含下划线（`q1\_core.py`），故一律经 `\detokenize` 处理，避免被当作数学下标。
% ★ v9：代码改用等宽字体 `\lstmonofam` —— `\ttfamily`（lmmono）**缺 ★／希腊字母／数学符号字形**，
%   代码注释里的这些符号会在 PDF 中**留空**（`Missing character` 警告）；改用覆盖更广的 DejaVu Sans Mono，
%   字体不可用时自动退回 `\ttfamily`（不影响编译）。**仅作用于附录 B 的代码块**。
\IfFontExistsTF{DejaVu Sans Mono}{\newfontfamily\lstmonofam{DejaVu Sans Mono}}{\let\lstmonofam\ttfamily}
\newcommand{\srcfile}[1]{%
  \edef\srcpath{\detokenize{#1}}%
  % ★ v11：代码字号由 `\footnotesize`（10pt）降为 `\scriptsize`（9pt）—— 实测 `code/*.py`
  %   第 4 行的 80 字符 `=` 横幅注释在 10pt 等宽下宽 481.6pt，超版心 455.2pt 达 **26.4pt**
  %   （日志 7 处 `Overfull \hbox`，附录 B 代码块内文字越出边框）；9pt 下 80 字符仅 433.4pt，
  %   且嵌入集内最长行恰为 80 字符 ⟹ 全部收回版心。`breaklines=true` 保留作二次保险。
  \IfFileExists{code/\srcpath}{\lstinputlisting[breaklines=true,basicstyle=\lstmonofam\scriptsize]{code/\srcpath}}%
  {\fbox{\parbox[c][1.1cm][c]{0.9\textwidth}{\centering\ttfamily\footnotesize 源程序 \detokenize{#1}（放入 code/ 目录后自动嵌入）}}}}

% ---- 图片：**依《A\_图表编号对照.md》的命名候选链自动解析** ----
% `\fgphchain{宽度}{候选1,候选2,候选3}`：依次尝试 `候选.pdf`／`候选.png`，命中即插入；
% 三个候选分别为 PDF 文件索引号（图NN，交付命名）／F 号／论文内索引号，皆缺则显示占位框
% （框内列出候选名，便于贵方对号投放）。**工程在任何情况下均可编译**。
\makeatletter
\newcommand{\fgphchain}[2]{%
  \def\fgfound{0}%
  \@for\fcand:=#2\do{%
    \ifnum\fgfound=0
      \IfFileExists{\fcand}{\includegraphics[width=#1]{\fcand}\def\fgfound{1}}{%
        \IfFileExists{\fcand.pdf}{\includegraphics[width=#1]{\fcand}\def\fgfound{1}}{%
          \IfFileExists{\fcand.png}{\includegraphics[width=#1]{\fcand}\def\fgfound{1}}{}}}%
    \fi}%
  \ifnum\fgfound=0
    \fbox{\parbox[c][2.6cm][c]{#1}{\centering 图占位\\\ttfamily\footnotesize #2}}%
  \fi}
\makeatother
\newcommand{\fgph}[2][0.82\textwidth]{\fgphchain{#1}{#2}}

% ---- 摘要页与关键词 ----
\newcommand{\keywords}[1]{\par\vspace{0.8em}\noindent\textbf{关键词：}#1\par}

% ---- 编号：**全部由 LaTeX 计数器在编译期动态生成**（源文件内不写死任何图/表/式号）----
%  章-序 连字符风格（式 (5-1)、图 5-1、表 5-1；附录自动成 表 A-1、图 C-1）；`\label`-`\ref` 自动跟随。
\numberwithin{equation}{section}
\numberwithin{figure}{section}
\numberwithin{table}{section}
% ---- ★ v7：题面强制表（表 1–表 6）的**独立计数器** ----
%     六张题面表由生成器在**组内**把 `\thetable` 改写为 `\arabic{reqtab}` ⟹ 全篇连续 1–6；
%     本文补充表（表 5-1）与附录表（表 A-1）仍走"章-序" ⟹ 编号纪律与《A\_图表编号对照.md》不变。
\newcounter{reqtab}
\renewcommand{\theequation}{\thesection-\arabic{equation}}
\renewcommand{\thefigure}{\thesection-\arabic{figure}}
\renewcommand{\thetable}{\thesection-\arabic{table}}

% ---- 表格：列宽**按各列内容加权分配**（旧实现各列均分 ⟹ 短列过宽、长列（说明／作用）过窄）----
% 用法：每张表前 `\settabw{2n}`（n ＝ 列数）先记下可用宽度，各列 p{} 再按月其**份额**取用。
\newlength{\tblw}
\newcommand{\settabw}[1]{\setlength{\tblw}{\dimexpr0.97\textwidth-#1\tabcolsep\relax}}

% ---- 列表（源程序语言高亮）----
\lstset{basicstyle=\ttfamily\footnotesize,breaklines=true,frame=single,
        numbers=left,numberstyle=\tiny,columns=fullflexible,keepspaces=true,
        % ★ v12：**消除 44 处 `Missing character`**（附录 B 代码注释的 ⟹×43 ＋ ✅×1 ——
        %   DejaVu Sans Mono 无这两个字形，PDF 中留空）。用 `literate` 在**排版层**换成
        %   数学／符号字体里的同义字形 ⟹ **不改动 `code/` 下任何源文件的字节**（与真源保持
        %   逐字节一致、哈希不变），仅影响渲染。
        literate={⟹}{{\ensuremath{\Longrightarrow}}}1 {✅}{{\ding{51}}}1}

% ---- 标题层级：章/节/子节（不加人为断页；表题在上、图题在下由生成器保证）----
% ★ v13：**标题字号/对齐按 kb 口径落地**（官方第八条"字号、字体、行距、颜色等不做统一要求"）
%   口径源（三处，已逐条读取原文）：
%   ①《必看！2026数学建模国赛获奖要点解析》要点48："一级标题黑体三号，二级标题黑体四号，三级标题黑体小四"
%   ②《2026国赛零基础获奖速成课讲义》表15/16："标题1 黑体三号(16pt) 居中、段前段后各12pt、加粗；
%      标题2 黑体四号(14pt) 左对齐、段前段后各6pt；标题3 黑体小四(12pt) 左对齐、段前段后各3pt"
%   ③《AI 论文全维度自查表》13.2："正文宋体小四，标题黑体，数字/公式 Times New Roman"
%   ⟹ 一级 三号16pt **居中**、二级 四号14pt、三级 小四12pt，均黑体（ctex 的 \bfseries → SimHei）。
\ctexset{section={format=\zihao{3}\bfseries\centering,aftername=\quad},
         subsection={format=\zihao{4}\bfseries,aftername=\quad},
         subsubsection={format=\zihao{-4}\bfseries,aftername=\quad}}
% ★ v13 归因：讲义表15/16 另有"段前段后各 12pt／6pt／3pt"一条，实测按此设置会使
%   两张长表页内分块超出 ⟹ 冒 2 处 `Overfull \vbox`（p9 142pt／p11 296pt）。
%   按用户既定优先级【**字体 ＞ 字号 ＞ 行间距**】，此条**属最低优先的间距项** ⟹ 回退为
%   ctex 默认段距，**保留**字体（黑体）与字号（三号／四号／小四）两项。**行距维持单倍**。
% ★ v13：**西文/数字改用 Times New Roman**（口径源 ③）；中文仍为 ctex 的 SimSun。
\setmainfont{Times New Roman}
\usepackage{newtxmath}
